\documentclass[11pt,a4paper]{article}
\usepackage[utf8]{inputenc}
\usepackage[T1]{fontenc}
\IfFileExists{lmodern.sty}{\usepackage{lmodern}}{}
\usepackage[spanish,es-tabla]{babel}
\usepackage{amsmath,amssymb}
\usepackage{siunitx}
\usepackage{graphicx}
\usepackage{booktabs}
\usepackage{caption}
\usepackage{float}
\usepackage{xcolor}
\usepackage[a4paper,margin=2.5cm]{geometry}
\usepackage[hidelinks]{hyperref}
\sisetup{output-decimal-marker={,},group-separator={\,},locale=DE}
\DeclareSIUnit\microstrain{\text{\ensuremath{\mu\varepsilon}}}
\newcommand{\pend}[1]{\textcolor{red}{\textbf{[PENDIENTE: #1]}}}
\graphicspath{{figures/}}

\title{Determinaci\'on del esfuerzo normal por flexi\'on en una viga instrumentada:\\ soluci\'on anal\'itica, experimental y por elementos finitos}
\author{Brayan \pend{apellidos} \and Samuel \pend{apellidos}\\[2pt]
Universidad Pontificia Bolivariana, Medell\'in\\ Elementos Finitos -- Prof.\ Elkin Alonso Taborda Espinosa}
\date{6 de octubre de 2026}

\begin{document}
\maketitle

\begin{abstract}
Se determina el esfuerzo normal por flexi\'on en la posici\'on de una galga extensiom\'etrica adherida a una viga tubular de aluminio 1100 en voladizo, sometida a una carga de \SI{49,77}{N}. El problema se resuelve de tres formas: con la teor\'ia de vigas de Euler--Bernoulli, con datos de laboratorio y con una simulaci\'on por elementos finitos (FEM). El valor anal\'itico es \SI{11,40}{MPa} (\SI{165,2}{\microstrain}); el experimental, \SI{12,55}{MPa}, difiere en \SI{+10,1}{\percent}. \pend{resultado FEM, error FEM y conclusi\'on tras completar la simulaci\'on}.
\end{abstract}

\section{Introducci\'on}
Una galga extensiom\'etrica mide la deformaci\'on unitaria longitudinal de la superficie sobre la que est\'a adherida, y el esfuerzo se obtiene con la ley de Hooke. En este trabajo se compara ese esfuerzo con el que predicen la mec\'anica de materiales y una simulaci\'on FEM, con el fin de validar los modelos y analizar las fuentes de error.

\section{Datos del montaje}
La viga es un tubo rectangular hueco de aluminio 1100, empotrado por medio de mordazas, con una pesa colgada de un tornillo cerca del extremo libre. La galga est\'a adherida en la cara superior. Las coordenadas se toman del modelo CAD (\texttt{VIGA\_FEM.step}), con el eje $y$ a lo largo de la viga.

\begin{table}[H]
\centering
\caption{Datos del montaje.}
\begin{tabular}{@{}lcc@{}}
\toprule
Par\'ametro & S\'imbolo & Valor \\
\midrule
M\'odulo de elasticidad & $E$ & \SI{69}{GPa} \\
Masa de la pesa & $m_p$ & \SI{5,000}{kg} \\
Masa del tornillo y arandelas & $m_a$ & \SI{0,073}{kg} \\
Gravedad & $g$ & \SI{9,81}{m/s^2} \\
Dimensiones exteriores de la secci\'on & $b_e\times h_e$ & \SI{70}{mm}$\times$\SI{30}{mm} \\
Espesor de pared & $t$ & \SI{1}{mm} \\
Longitud total & $L$ & \SI{746}{mm} \\
Borde libre de la mordaza & $y_m$ & \SI{141,83}{mm} \\
Centro de la galga (marca de 198 a \SI{213}{mm}) & $y_g$ & \SI{205,5}{mm} \\
Eje del tornillo (punto de carga) & $y_P$ & \SI{711}{mm} \\
Rejilla sensible & -- & \SI{2,5}{mm}$\times$\SI{10}{mm} \\
\bottomrule
\end{tabular}
\end{table}

\section{Parte 1: modelo anal\'itico}
\subsection{Hip\'otesis}
Se usa la teor\'ia de Euler--Bernoulli con material el\'astico lineal, peque\~nas deformaciones y secci\'on prism\'atica. La carga es est\'atica y vertical. El peso propio se excluye porque la galga se puso en cero con la viga ya montada, de modo que el ensayo mide \'unicamente el incremento debido a la carga. En la cara superior, que es una superficie libre, el esfuerzo es uniaxial ($\sigma_x$), por lo que el esfuerzo de Von Mises coincide con $|\sigma_x|$.

\subsection{Carga}
\begin{align}
m &= m_p + m_a = \SI{5,073}{kg},\\
P &= m\,g = 5{,}073\times 9{,}81 = \SI{49,766}{N}.
\end{align}
La pesa cuelga de dos cuerdas sim\'etricas inclinadas $20^\circ$ respecto a la vertical (figura~\ref{fig:cargas}). La fuerza en cada cuerda es
\begin{equation}
F=\frac{P}{2\cos 20^\circ}=\SI{26,48}{N},
\end{equation}
sus componentes horizontales se anulan y la suma de las verticales es exactamente $P$.

\begin{figure}[H]
\centering
\includegraphics[width=0.35\textwidth]{fig_cargas_fusion.png}
\caption{Cargas aplicadas en el modelo: dos fuerzas que parten de los extremos del tornillo y convergen hacia el nudo de la pesa.}
\label{fig:cargas}
\end{figure}

\subsection{Secci\'on transversal}
La secci\'on es un rect\'angulo hueco de \SI{70}{mm}$\times$\SI{30}{mm} con espesor \SI{1}{mm}. La viga se ensaya con la altura de \SI{30}{mm} en vertical, as\'i que la flexi\'on ocurre alrededor del eje de menor inercia:
\begin{align}
I &= \frac{b_e h_e^3-(b_e-2t)(h_e-2t)^3}{12}=\frac{70\cdot 30^3-68\cdot 28^3}{12}=\SI{33105,3}{mm^4},\\
c &= \frac{h_e}{2}=\SI{15}{mm}.
\end{align}

\subsection{Momento flector, esfuerzo y deformaci\'on}
El momento en la galga depende de su distancia al punto de carga:
\begin{align}
a &= y_P-y_g = 711-205{,}5=\SI{505,5}{mm},\\
M(y_g) &= P\,a = 49{,}766\times 505{,}5=\SI{25156,8}{N.mm},\\
\sigma &= \frac{M\,c}{I}=\frac{25156{,}8\times 15}{33105{,}3}=\SI{11,399}{MPa},\\
\varepsilon &= \frac{\sigma}{E}=\frac{11{,}399}{69000}=\SI{165,2}{\microstrain}.
\end{align}
Como el momento es lineal en $y$, el promedio sobre la rejilla coincide con el valor en su centro; en los extremos de la marca de galga ($y=198$ y \SI{213}{mm}) el esfuerzo es \SI{11,57}{MPa} y \SI{11,23}{MPa}.

\subsection{Diagramas de fuerzas internas}
El cortante es constante e igual a $P$ entre la mordaza y la carga, y el momento flector decrece linealmente desde $M_0=\SI{28,33}{N.m}$ en el borde de la mordaza hasta cero en el punto de carga (figura~\ref{fig:vm}). La galga se ubica en una zona de momento alto, donde $M=\SI{25,16}{N.m}$.

\begin{figure}[H]
\centering
\includegraphics[width=0.8\textwidth]{fig_diagramas_V_M.png}
\caption{Esquema de la viga y diagramas de cortante $V(y)$ y momento flector $M(y)$.}
\label{fig:vm}
\end{figure}

\section{Parte 2: resultados experimentales}
El archivo del ensayo contiene 7\,744 muestras a \SI{100}{Hz} (\SI{77}{s}). La se\~nal muestra una l\'inea base, una rampa de carga, una meseta estable y la descarga (figura~\ref{fig:exp}). Se promedi\'o el strain en la meseta ($t=24$ a \SI{53}{s}, $n=2\,901$ muestras) y se le rest\'o la l\'inea base inmediatamente anterior a la carga ($t=11$ a \SI{16,5}{s}), porque la se\~nal deriva de \SI{0,4}{\microstrain} al inicio a \SI{3,5}{\microstrain} antes de cargar.

\begin{table}[H]
\centering
\caption{Resultados experimentales.}
\begin{tabular}{@{}lcc@{}}
\toprule
Magnitud & Sin l\'inea base & Con l\'inea base \\
\midrule
Strain promedio en la meseta & \SI{185,31}{\microstrain} & \SI{181,85}{\microstrain} \\
Desviaci\'on est\'andar en la meseta & \multicolumn{2}{c}{\SI{0,86}{\microstrain}} \\
Esfuerzo $\sigma=E\varepsilon$ & \SI{12,79}{MPa} & \SI{12,55}{MPa} \\
\bottomrule
\end{tabular}
\end{table}

\begin{figure}[H]
\centering
\includegraphics[width=0.8\textwidth]{fig_experimental.png}
\caption{Deformaci\'on medida por la galga, con la ventana de promediado y la l\'inea base.}
\label{fig:exp}
\end{figure}

\section{Parte 3: simulaci\'on por elementos finitos}
\subsection{Modelo CAD}
El modelo se construy\'o en Autodesk Fusion a partir de la geometr\'ia de la viga (tres s\'olidos contiguos), el tornillo, dos arandelas y una tuerca de cierre idealizada. La marca de galga de \SI{5}{mm}$\times$\SI{15}{mm} ocupa $y=198$ a \SI{213}{mm} en la cara superior, y el agujero de \SI{8}{mm} est\'a en $y=\SI{711}{mm}$. \pend{im\'agenes del modelo CAD desde la orientaci\'on del ensayo}

\subsection{Condiciones de la simulaci\'on}
Est\'atico lineal con aluminio 1100 y $E=\SI{69}{GPa}$; apoyo fijo en la zona de mordaza; dos fuerzas remotas de \SI{26,48}{N} a $20^\circ$ de la vertical aplicadas en $y=\SI{711}{mm}$, sin gravedad (coherente con el tare de la galga). \pend{confirmar apoyo, tipo de contactos y elemento}

\subsection{Independencia de malla}
\begin{table}[H]
\centering
\caption{Estudio de convergencia (esfuerzo en el centro de la galga).}
\begin{tabular}{@{}cccc@{}}
\toprule
Tama\~no de elemento [mm] & Nodos & $\sigma_{VM}$ [MPa] & Variaci\'on [\%] \\
\midrule
25 & \pend{} & \pend{} & -- \\
20 & \pend{} & \pend{} & \pend{} \\
15 & \pend{} & \pend{} & \pend{} \\
10 & \pend{} & \pend{} & \pend{} \\
5 & \pend{} & \pend{} & \pend{} \\
\bottomrule
\end{tabular}
\end{table}
\pend{gr\'afica de convergencia y valor final en la galga}

\section{Parte 4: comparaci\'on y an\'alisis}
\begin{table}[H]
\centering
\caption{Comparaci\'on de resultados (error respecto al anal\'itico).}
\begin{tabular}{@{}lcc@{}}
\toprule
M\'etodo & $\sigma$ [MPa] & Error [\%] \\
\midrule
Anal\'itico & \SI{11,40}{} & -- \\
Experimental (con l\'inea base) & \SI{12,55}{} & $+10{,}1$ \\
Experimental (sin l\'inea base) & \SI{12,79}{} & $+12{,}2$ \\
FEM & \pend{} & \pend{} \\
\bottomrule
\end{tabular}
\end{table}

\subsection{Causas de las diferencias}
\textbf{Espesor de pared.} La inercia depende casi linealmente del espesor: una variaci\'on de \SI{-0,05}{mm} eleva el esfuerzo te\'orico \SI{4,8}{\percent}. Un espesor real de \SI{0,90}{mm} bastar\'ia para que el valor anal\'itico llegara a \SI{12,56}{MPa}, casi id\'entico al experimental; conviene medirlo con calibrador. Los radios de esquina del perfil, ignorados en el modelo, tambi\'en reducen $I$.\\
\textbf{Galga y adquisici\'on.} El factor de galga configurado, la orientaci\'on de la galga, la resistencia de los cables y la deriva de la se\~nal (\SI{3,5}{\microstrain} antes de cargar y \SI{1,3}{\microstrain} despu\'es de descargar) aportan entre 1 y $\SI{2}{\percent}$.\\
\textbf{Geometr\'ia y carga.} Un error de \SI{3}{mm} en el brazo cambia el resultado solo \SI{0,6}{\percent}, y la incertidumbre de la masa (aprox.\ \SI{1}{\percent}) es menor que la diferencia observada.\\
\textbf{Peso propio.} Est\'a excluido del ensayo (tare); en el FEM debe apagarse la gravedad, ya que aportar\'ia \SI{0,34}{MPa} (aprox.\ \SI{3}{\percent}).\\
\pend{causas espec\'ificas del FEM: tipo de elemento, malla en el espesor de \SI{1}{mm}, condici\'on de apoyo y contactos del tornillo}

\section{Conclusiones}
\pend{redactar cuando est\'e el resultado FEM}

\appendix
\section{Anexos}
Archivos con los datos (\texttt{Viga4SamuelBrayan.csv}), modelo CAD en formato original (\texttt{.f3d}) y universal (\texttt{.step}), y archivos de la simulaci\'on FEM. Repositorio de trabajo: \url{https://github.com/brayansan00/galgalab}.

\end{document}
